
\medskip

%\emph{Cet exercice est un questionnaire à choix multiple.\\
%Pour chaque question, une seule des quatre réponses proposées est exacte.\\
%Le candidat indiquera sur sa copie le numéro de la question et la réponse choisie.\\
%Aucune justification n'est demandée. \\
%Une réponse fausse, une réponse multiple ou l'absence de réponse à une question ne rapporte ni n'enlève de point.}
%
%\emph{Les cinq questions sont indépendantes.}
%
%\medskip
%
%L'espace est rapporté à un repère orthonormé \Oijk.
%
%\medskip

\begin{enumerate}
\item On considère les points A(1~;~0~;~3) et B(4~;~1~;~0).

Une représentation paramétrique de la droite (AB) est:

%\begin{center}
%\begin{tabularx}{\linewidth}{*{2}{X}}
%\textbf{a.~}$\left\{\begin{array}{l c l}
%x &=& \phantom{-}3 + t \\
%y &=& \phantom{-}1\\
%z &=& - 3 +3t
%\end{array}\right.$ avec $t \in \R$&\textbf{b.~}$\left\{\begin{array}{l c l}
%x &=& 1+4t\\
%y &=&\phantom{1 + 4}t\\
%z &=& 3\end{array}\right.$ avec $t \in \R$\\
%\textbf{c.~}$\left\{\begin{array}{l c l}
%x &=& 1 + 3t\\
%y &=&\phantom{1 + 3}t\\
%z&=&3 - 3t
%\end{array}\right.$ avec $t \in \R$&\textbf{d.~}$\left\{\begin{array}{l c l}
%x&=&4 + t\\
%y&=&1\\
%z&=&3 - 3t
%\end{array}\right.$ avec $t \in \R$
%\end{tabularx}
%\end{center}

On a $\vect{\text{AB}}\begin{pmatrix}3\\1\\-3\end{pmatrix}$.
$M(x~;~y~;~z) \in (\text{AB}) \iff \vect{\text{A}M} = t\vect{\text{AB}}, \: t \in \R \iff $

$\left\{\begin{array}{l c l}x - 1&=&3t\\y - 0&=&1t\\z - 3&=&- 3t, \end{array}\right. \: t \in \R\iff \left\{\begin{array}{l c l}
x &=& 1 + 3t\\
y &=&\phantom{1 + 3}t\\
z&=&3 - 3t
\end{array}\right.\: t \in \R$ Réponse \textbf{c.~}

On considère la droite $(d)$ de représentation paramétrique 

$\left\{\begin{array}{l c l}
x&=&3 + 4t\\
y&=&\phantom{3 + }6t\\
z&=&4 - 2 t
\end{array}\right.$ avec $t \in \R$

\item %Parmi les points suivants, lequel appartient à la droite $(d)$ ?
%\begin{center}
%\begin{tabularx}{\linewidth}{*{4}{X}}
%\textbf{a.~} M(7~;~6~;~6) &\textbf{b.~} N(3~;~6~;~4) &\textbf{c.~}P$(4~;~6~;~-2)$&\textbf{d.~} R$(-3~;~-9~;~7)$
%\end{tabularx}
%\end{center}
Réponse \textbf{d}.
\item %On considère la droite $(d')$ de représentation paramétrique 

%$\left\{\begin{array}{l c l}
%x&=&-  2 + 3k\\y&=&- 1 - 2k\\
%z&=&\phantom{-}1 + k
%\end{array}\right.$ avec $k \in \R$
%
%Les droites $(d)$ et $(d')$ sont:
%
%\begin{center}
%\begin{tabularx}{\linewidth}{*{4}{X}}
%\textbf{a.~} sécantes &\textbf{b.~}  non coplanaires&\textbf{c.~}  parallèles&\textbf{d.~}  confondues
%\end{tabularx}
%\end{center}
$\bullet~~$Les deux droites ont pour vecteurs directeurs respectifs : $\begin{pmatrix}4\\6\\-2\end{pmatrix}$ et $\begin{pmatrix}3\\-2\\1\end{pmatrix}$ : ces vecteurs ne sont pas colinéaires, les droites ne sont pas parallèles.

$\bullet~~$Les deux droites sont sécantes s'il existe deux réels $t$ et $k$ tels que :

$\left\{\begin{array}{l c l}
3 + 4t&=&- 2 + 3k\\
6t&=&- 1 - 2k\\
4 - 2t&=&1 + k
\end{array}\right. \iff \left\{\begin{array}{l c l}
4t&=&3k - 5\\6t&=&- 2k - 1\\
- 2t&=&k - 3\end{array}\right.$

Par différence de la ligne 1 et de la ligne 2, on obtient$ -2t = 5k - 4 = k - 3$ (ligne 3) soit $4k = 1 \iff k = \dfrac14$, puis comme $- 2t = k - 3 = \dfrac14 - 3 = - \dfrac{11}{4}$, d'où $t = \dfrac{11}{8}$, mais la première ligne donne $4t = 3k - 5 = \dfrac34 - 5 = - \dfrac{17}{4}$, d'où $t = \dfrac{17}{8}$ : ceci n'est pas possible : les deux droites ne sont pas sécantes.

$\bullet~~$Les deux droites ne sont pas confondues : il reste la réponse \textbf{b.}

\item  %On considère le plan $(P)$ passant par le point I(2~;~1~;~0) et perpendiculaire à la droite $(d)$.

%Une équation du plan $(P)$ est:
%
%\begin{center}
%\begin{tabularx}{\linewidth}{*{4}{X}}
%\textbf{a.~} $2x + 3y - z - 7=0$&\textbf{a.~} $- x + y - 4z + 1 = 0$\\
%\textbf{b.~} $4x + 6y - 2z + 9 = 0$&\textbf{d.~} $2x + y + 1 = 0$.
Le vecteur $\vect{n}$ vecteur directeur de $(d)$ est un vecteur normal au plan $(P)$ . On a donc 

$M(x~;~y~;~z) \in (P) \iff \vect{\text{I}M} \cdot \vect{n} = 0 \iff 4(x - 2) + 6(y - 1) - 2(z - 0) = 0 \iff 4x - 8 + 6y - 6 - 2z = 0 \iff 4x + 6y -2z - 14 = 0 \iff 2x + 3y - z - 7 = 0$ : réponse \textbf{a.}
\end{enumerate}

\bigskip

